
\subsubsection{Métodos de función de penalización} 
Tal y como se explicó en el caso de optimización de un único nivel con restricciones, los métodos de penalización definen un nuevo problema en el que el objetivo original se combina con un nuevo término, que siempre es mayor o igual que cero y que está relacionado con la restricción. Si la restricción se satisface, la penalización vale cero. En caso contrario la penalización se activa, tomado mayores valores cuanto mayor sea el nivel de violación de las restricciones. Para el caso de la BLO, se considera un escenario sin restricciones y se denota como $\gamma$ al coeficiente que refleja el peso de la penalización. El problema penalizado de un solo nivel en dicho escenario es
\begin{align}\label{eq:blo_pen_introblo}
    \min_{x,y} f(x,y) + \gamma P(x,y),
\end{align}
donde la función de penalización $P(x,y)$ está asociada al problema del nivel inferior. En la literatura existen varias opciones para definir esta función de penalización \cite{aiyoshi1981hierarchical,aiyoshi1984solution,ishizuka1992double}. Para el caso concreto en el que el problema inferior no tiene restricciones, una función de penalización común es $p(x,y)=\|\nabla_y g(x,y)\|^2$. De este modo, si el valor de $\gamma$ es muy elevado, la solución de \eqref{eq:blo_pen_introblo} tenderá a hacer que $\nabla_y g(x,y)\rightarrow 0$, garantizando que se satisfagan las condiciones de optimalidad del nivel inferior. Naturalmente, la norma 2 en la expresión anterior puede sustituirse por cualquier otra norma. Asimismo, si el problema en el nivel inferior tiene restricciones, la penalización se aplicará sobre un vector que contenga el gradiente del lagrangiano del nivel inferior, así como el resto de las condiciones KKT. En el contexto de métodos de penalización para BLO, también es común definir la función de penalización como $p(x,y)=g(x,y)-v(x)$, donde $v(x)$ es una aproximación de $\min_y g(x,y)$ conocida como \emph{función de valor}. De esta forma, $p(x,y)=0$ sólo si $y=\arg\min_z g(x,z)$, siendo mayor que cero en cualquier otro caso. Al igual que cuando se penaliza la norma del gradiente, existen algoritmos que modifican la función de penalización utilizada (por ejemplo, aplicando una transformación no lineal y creciente a $g(x,y)-v(x)$), de forma que se pueda modular la sensibilidad del algoritmo a la violación de la condición de optimalidad. El algoritmo diseñado en este TFG pertenece a la familia de métodos penalizados utilizando la función de valor, por lo que varios de estos aspectos se discutirán en más detalle en el capítulo \ref{chap:BLOCC}.


\section{Revisión histórica del estado del arte}\label{sec:evol_hist_fundBLO}


Una vez descritos en detalle los elementos que componen un problema de BLO y las técnicas más habituales para abordarlos, se procede a realizar una revisión histórica del estado del arte. Este análisis proporciona un contexto sobre la evolución y el estado actual de los métodos de solución para problemas de BLO.

La primera formulación de un problema de optimización binivel se remonta a más de 50 años atrás \cite{bracken1973mathematical}. Sin embargo, los trabajos sobre juegos líder-seguidor de Von Stackelberg \cite{von2010market,stackelberg1952theory} ya identifican los principales elementos de los problemas BLO. Una discusión temprana sobre problemas binivel se encuentra en \cite{candler1977multilevel}, donde se reconoce la dificultad intrínseca de estos problemas y sus diferencias con los problemas de un solo nivel. Los autores de \cite{candler1977multilevel} observaron que incluso con variables continuas y funciones objetivo y restricciones lineales, el conjunto factible de los problemas binivel puede ser no convexo y disconexo. De hecho, los resultados formales de complejidad indican que esta clase simplificada de problemas binivel es fuertemente NP\footnote{NP: \emph{Nondeterministic polynomial time.}}-difícil . En \cite{candler1977multilevel} también se propuso uno de los primeros algoritmos para resolver problemas BLO lineales, similar al método símplex. Estudios posteriores, como \cite{bialas1978multilevel}, profundizaron en este tipo de algoritmos.

En 1981, \cite{fortuny1981representation} diseñó un nuevo método para problemas binivel convexos-cuadráticos. La estructura cuadrática del nivel inferior da lugar a condiciones KKT lineales, habilitando la sustitución del problema de optimización inferior por un sistema de ecuaciones lineales. No obstante, dado que solo se deben considerar las restricciones activas, es necesario introducir variables binarias para seleccionar las restricciones activas. De esta forma, el método de \cite{fortuny1981representation} resulta en un problema MIP de un solo nivel que puede ser abordado por \emph{solvers} estándar. Otros trabajos, como \cite{bard1990branch,bard1988convex,edmunds1991algorithms,hansen1992new}, profundizaron en esta metodología, siendo aún hoy en día una estrategia popular para resolver problemas de BLO convexos en el nivel superior y cuadráticos en el nivel inferior. Posteriormente, \cite{ye1995optimality,dempe1992necessary,dempe1992optimality,yezza1996first,outrata1993necessary} desarrollaron las condiciones de optimalidad de primer y segundo orden para programas binivel que involucran problemas sin restricciones convexos-cuadráticos.

Además de los enfoques basados en KKT, se han propuesto varios métodos de descenso para resolver problemas de optimización binivel. Los trabajos en \cite{kolstad1990derivative,savard1990steepest} adaptaron la noción de descenso más pronunciado al caso de los programas binivel. Recientemente, los enfoques basados en descenso por gradiente han sido objeto de estudio. Una línea de trabajo se centra en utilizar variantes del método de IGD, donde la idea clave es calcular (aproximar) el gradiente del problema BLO (referido en la literatura como hipergradiente) mediante el teorema de la función implícita --véase \cite{pedregosa2016hyperparameter}, así como la discusión en torno a la ecuación \eqref{E:IGD_resultante} proporcionada en esta memoria--, evitando el cálculo explícito de la solución del nivel inferior y facilitando por tanto la optimización del problema del nivel superior.
IGD se ha utilizado principalmente para problemas BLO sin restricciones \cite{hong2022twotimescale,ghadimi2018approximation,ji2020provably,chen2022single,chen2021closing}, aunque también existen trabajos para resolver problemas de BLO con restricciones desacopladas \cite{khanduri2023linearly,kwon2023penalty,shen2023penaltybased}. %, e incluso acopladas \cite{xiao2022alternating,xu2023efficient}, pese a tener múltiples limitaciones. 
En cuanto al tiempo de cómputo, la convergencia en tiempo finito de los métodos de IGD se estableció por primera vez en \cite{ghadimi2018approximation} para problemas de nivel inferior fuertemente convexos y sin restricciones. Trabajos posteriores, como \cite{ji2020provably,chen2022single,chen2021closing,hong2020ac,khanduri2021near,shen2022single,Li2022fully,sow2022convergence,yang2023accelerating,ji2021bilevel}, mejoraron las tasas de convergencia, en gran parte de los casos, realizando suposiciones adicionales con respecto a la estructura y propiedades del problema.

%Otro grupo de trabajos se basa en la diferenciación iterativa (ITD) \cite{maclaurin2015gradient, franceschi2017forward, nichol2018firstorder, shaban2019truncated}, que estima el hipergradiente diferenciando todo el algoritmo iterativo utilizado para resolver el problema de nivel inferior con respecto a las variables de nivel superior. La garantía de tiempo finito se estableció por primera vez en \cite{grazzi2020iteration,liu2021towards,ji2022will,bolte2022automatic}.

En cuanto a enfoques penalizados, los primeros intentos se produjeron en los años 80 en los estudios realizados en \cite{aiyoshi1981hierarchical,aiyoshi1984solution}. Estos trabajos reemplazan el problema del nivel inferior por un término de penalización en la función objetivo del nivel superior, obteniendo un problema de optimización de un solo nivel \cite{ye1997exact,liu2021value,mehra2021penalty,liu2022bome,kwon2023fully}. El término de penalización se introduce con el objetivo de asegurar la optimalidad del problema del nivel inferior. Otros estudios que se basan en el uso de funciones de penalización para problemas BLO lineales son \cite{xu2012nonlinear,maldonado2016analyzing,shen2023penaltybased,gao2022value,lu2024slm}.


%Más tarde, el método de doble penalización en \cite{ishizuka1992double} penalizó el objetivo de ambos niveles y luego sustituyó la reformulación del problema del nivel inferior por sus condiciones KKT. 

% En \cite{xu2012nonlinear}, los autores convierten el problema de BLO lineal en un problema de BLO penalizado, y luego resuelven una serie de problemas bilineales para encontrar el óptimo. En \cite{maldonado2016analyzing}, los autores reducen el programa binivel lineal a un solo nivel utilizando las condiciones KKT, y luego añaden la condición de holgura complementaria a la función objetivo de nivel superior con una penalización. El problema penalizado se maneja luego utilizando una serie de programas lineales. Los métodos basados en penalizaciones también han surgido como un enfoque prometedor para la BLO. 
 
%La línea de trabajo iniciada en \cite{ye1997exact} y  continuada en \cite{liu2021value,mehra2021penalty,liu2022bome,kwon2023fully,shen2023penaltybased,gao2022value,lu2024slm} reformuló el BLO original como distintos problemas de un solo nivel (cada uno con diversos términos de penalización) y empleó métodos de primer orden para resolverlos. 

Aunque se han logrado avances significativos para problemas de BLO sin restricciones, el análisis para BLO con restricciones es más limitado. Las restricciones de nivel superior de la forma~\eqref{E:upper_cons} fueron consideradas en \cite{chen2022single,hong2020ac,chen2022fast}. Para BLO con restricciones desacopladas en el nivel inferior, se han propuesto algoritmos en \cite{khanduri2023linearly,kwon2023penalty,shen2023penaltybased}. 
%Para las restricciones desacopladas de nivel inferior, el algoritmo propuesto en \cite{khanduri2023linearly} logró una convergencia asintótica, mientras que \cite{kwon2023penalty, shen2023penaltybased} emplearon una reformulación basada en penalizaciones y consideraron tanto restricciones desacopladas de nivel superior como inferior. 
%
Sin embargo, la literatura sobre BLO con CCs es escasa, a pesar de que estos problemas modelan muchos escenarios reales. En los últimos años, se ha suscitado un interés creciente en estos problemas, como se observa en estudios recientes \cite{liu2023value,xiao2022alternating,yao2024constrained,yao2024overcoming}. Los detalles de BLO con CCs se profundizarán en la sección introductoria del capítulo \ref{chap:BLOCC}.

%BVFSM \cite{liu2023value} utilizó un método basado en penalizaciones para evitar el cálculo del Hessiano, como lo hacen los métodos IGD. Sin embargo, solo se logró una convergencia asintótica. GAM \cite{xu2023efficient} investigó el método IGD bajo restricciones de desigualdad, pero tampoco logró proporcionar resultados de convergencia en tiempo finito. AiPOD \cite{xiao2022alternating} también aplicó IGD y logró con éxito una convergencia en tiempo finito, pero solo consideró restricciones de igualdad. LV-HBA \cite{yao2024constrained} consideró restricciones de desigualdad y construyó una reformulación basada en penalizaciones. No obstante, empleó una proyección conjunta de las variables de los niveles superior sobre las restricciones acopladas, lo cual es computacionalmente ineficiente cuando hay muchas restricciones y/o variables o cuando las restricciones no son convexas.
%\cite{yao2024overcoming} utilizó la teoría de dualidad de Lagrange para construir un nuevo término de penalización suavizado. También evitó la proyección conjunta sobre las restricciones acopladas y logró la misma tasa de convergencia que LV-HBA. Sin embargo, no logra cuantificar la relación entre la relajación de esta penalización y la relajación de la optimalidad de nivel inferior, y no garantiza la factibilidad del nivel inferior.

Paralelamente, se han desarrollado algoritmos evolutivos para abordar problemas de optimización con una estructura matemática desfavorable debido a su capacidad para afrontar problemas complejos. Algunos trabajos se basan en métodos de enjambre \cite{sinha2014bilevel}, mientras que otros utilizan metamodelos para resolver una versión simplificada y tratable del problema original \cite{wang2005evolutionary}.

La revisión realizada es, inevitablemente, limitada y sesgada por la temática de este TFG. Los lectores interesados en una descripción más completa del estado del arte de la BLO son remitidos a los tutoriales y revisiones realizados en \cite{sinha2018review,colson2007overview,bard2013practical}.